\chapter{Childcare, Eldercare, and Presence}
\label{ch:14}

Care for those who cannot care for themselves, whether because they are young, old, ill, or disabled, is the domain in which the temptation to delegate is strongest and the grounds for restraint are most evident. It is strongest because the demand is large and growing: populations are aging in many regions, the number of persons requiring assistance with the activities of daily life is rising, and the supply of paid and unpaid carers is not rising with it. It is a domain of restraint because what is being provided is not only a service but a relation, and the relation is not equivalent to the service. The distinction is easily stated and difficult to maintain under pressure. A person who is bathed, fed, moved, and medicated has received services. A person who is bathed, fed, moved, and medicated by someone who knows them, notices a change in their expression, converses with them, and regards them as a person has received care. The principle of the book places the first within the range of what may be assisted by machines and the second beyond the range of what may be replaced by them.

To ground the principle in the details of practice, it is useful to distinguish the tasks involved in care by their relation to the three categories introduced in Chapter 13. Some tasks are physical and well defined, such as lifting a person from a bed to a chair, which poses a substantial risk of injury to the carer. Assistance with such tasks by a machine that is able to bear load and that limits its own force is consistent with the principle, and is among its clearest applications, since the strain of lifting is a leading cause of injury and departure from the profession. Some tasks are monitoring tasks, such as watching for falls, detecting wandering, or recording vital signs, and sensors with software that flags anomalies can extend the reach of a carer without replacing the carer's response. Some are tasks of companionship and conversation, in which the matter is delicate: an automated conversational partner may relieve loneliness in some cases, and there is evidence of benefit in particular settings, but the substitution of such a partner for human contact is a different thing from its supplementation, and the temptation of the institution that is short of staff to treat the one as the other must be explicitly resisted. And some tasks are decisions about the person's life: where they will live, what treatment they will accept, whom they will see, how their money is used. These are not delegable to a system, and the more vulnerable the person the more care is owed to the process by which they are made and to the person's participation in it.

Children raise the same questions with additional urgency, since early relations form the basis of later capacities. The research literature on early development emphasizes the importance of responsive caregiving, in which an adult attends to the child's signals and responds contingently and consistently, and the absence of such care in early life has lasting effects. Devices that entertain or soothe children may be useful as aids, but they do not provide contingent response in the relevant sense, and a household or institution that relies on them for care in the absence of adults is in effect providing less care. The program accordingly sets a floor for the staffing of early childhood settings in terms of the number of adults per child, independent of any technology the setting adopts, and it regards technology as legitimate to the extent that it relieves adults of tasks other than attending to children. The same reasoning applies to the use of automated systems in the supervision of children online, where filtering and monitoring tools may reduce exposure to some harms but cannot replace the attention of a parent or teacher who knows the child.

A technical matter of some importance arises when care is provided at a distance. The architecture of Chapter 13 contemplates a human operator who directs a robotic embodiment remotely, a configuration that makes it possible for a skilled carer in one place to assist a person in another, and for a specialist to examine and treat a patient far from a hospital. The advantages are real, in particular for rural and underserved regions. The limits, however, are partly physical and partly physiological. Any communication channel has a delay, and a feedback loop with a delay can become unstable if the gain is too high: an operator who responds to what the remote embodiment appears to be doing, when it is actually doing something different by the time the response arrives, is liable to overcorrect, and the oscillation may grow. The formalization states the criterion for a simple loop. The practical consequence is that the functions requiring fast response, such as the limitation of force on contact with a person, the maintenance of balance, and the stopping of motion in the event of loss of signal, must be performed by controllers local to the embodiment, which run at rates far above those at which a human can respond, and the human operator must be confined to the slower, higher-level choices of what to do and when. This is the technical expression of the allocation of authority in the architecture, and it is a case where the engineering constraint and the ethical principle point in the same direction.

The sustainability of care work also requires attention, because the demand for care will not be met by the replacement of carers. It will be met, if at all, by changes in the conditions that determine who becomes a carer and how long they remain. Care work is at present poorly paid, physically demanding, emotionally exhausting, and weakly protected in many jurisdictions, and a substantial share of it is performed without pay by family members, disproportionately women. The introduction of machines that reduce the physical burden and administrative load can improve the working conditions of carers and make the work more attractive, but there is no guarantee that productivity gains will be used for this purpose and not to reduce staffing. The use of the gains is a matter of institutional design, and it is governed by the same consideration as in the earlier chapters: who decides, and who is accountable. A public body that purchases care technology should attach conditions about staffing levels and working conditions to the purchase and should publish the outcomes, so that the effects on care recipients and carers alike are visible.

Finally, the protection of dignity and privacy in care settings sets requirements that technology can either meet or violate. Monitoring technologies that observe a person in the bedroom or bathroom create a record of intimate life, and the person's consent to the monitoring is meaningful only if the person is capable of giving it and has a real alternative. The default for persons who lack capacity should be the least intrusive means consistent with safety, applied with the involvement of a guardian or advocate, with regular review. Data from care settings should be held under the strict limits applicable to health information, should not be used for purposes other than care without consent, and should not be sold. A technology that cannot operate within these limits does not belong in the setting, whatever its advantages in efficiency.

\section*{Formalization}

Consider a remote loop in which the operator's command acts on the embodiment through a channel with round-trip delay $\tau>0$. The simplest model of the resulting position error $x(t)$ under proportional correction with gain $K>0$ is the delay differential equation
\[
\dot x(t)=-K\,x(t-\tau).
\]

\begin{proposition}
The equation is asymptotically stable if and only if $K\tau<\pi/2$. For $K\tau=\pi/2$ it has a sustained oscillation of angular frequency $K$, and for $K\tau>\pi/2$ it has oscillating solutions of growing amplitude. The largest stable gain is $K_{\max}=\pi/(2\tau)$.
\end{proposition}

\begin{proof}
Substituting $x=e^{st}$ gives the characteristic equation $s+Ke^{-s\tau}=0$. At the stability boundary $s=i\omega$ with $\omega>0$ (the root $s=0$ is excluded since $K>0$). Then $i\omega+K(\cos\omega\tau-i\sin\omega\tau)=0$, whose real and imaginary parts are $K\cos\omega\tau=0$ and $\omega=K\sin\omega\tau$. The first gives $\omega\tau=\pi/2+n\pi$; the second requires $\sin\omega\tau>0$, so $n$ is even and the smallest root is $\omega\tau=\pi/2$ with $\omega=K$, so $K\tau=\pi/2$. For $K\tau<\pi/2$ no root has nonnegative real part: at $K=0^+$ the single root is $s\approx-K<0$ and the roots in the infinite chain originate at $\mathrm{Re}\,s=-\infty$; roots can cross the imaginary axis only at $K\tau=\pi/2+2n\pi$, the smallest of which is $\pi/2$. For $K\tau>\pi/2$ the crossing is from left to right, since $\mathrm{d}\,\mathrm{Re}\,s/\mathrm{d}K>0$ at the crossing, which gives growing oscillations.
\end{proof}

The consequences are quantitative. With round-trip delay $\tau=0.2$ s, which is typical of intercontinental links with processing and display, $K_{\max}=\pi/0.4\approx7.85\ \mathrm{s^{-1}}$, which corresponds to a correction time constant not shorter than about $0.13$ s. With delay $\tau=0.05$ s the maximum gain rises to about $31.4\ \mathrm{s^{-1}}$. A contact-force limiter, which must respond within milliseconds to prevent injury, would require $K$ of the order of $10^2$--$10^3\ \mathrm{s^{-1}}$, and by the proposition it cannot be placed in a loop that includes a human operator at any plausible delay. The division of control between local reflex layers and remote supervision is thus determined by $K\tau<\pi/2$ together with the response time that safety demands, and does not depend on an appeal to values.

A second quantity concerns staffing. Let $N$ persons need care, each requiring $c$ hours of attendance per week, of which a fraction $\theta\in[0,1]$ is attendance of the relational kind that is not delegated. The weekly human hours required are at least $N\theta c$, and if the proportion of carers' working hours spent on non-relational tasks is reduced by a labor-saving device that reduces these tasks by a fraction $\beta$ of the total, the hours released are $\beta Nc$ and may be reallocated to relational time only if $\beta Nc$ is allowed by the institution to be so used. In an institution with fixed staffing, relational time per person rises from $\theta c$ to $\theta c+\beta c$ if the released time is reallocated, and remains $\theta c$ if the staffing is reduced by $\beta Nc/H$, with $H$ the hours per carer per week. The difference between the two outcomes is the entire benefit of the technology to the persons cared for, which is the quantitative content of the requirement that purchase conditions secure the use of the gain.
